\chapter{Obsolete Straight-Coordinate Generalized Gradients}
\label{s:gen.grad.old}

This appendix documents the obsolete straight-coordinate generalized-gradient (GG) formalism of
Venturini and Dragt\cite{b:gen.grad}, which was used by the removed \vn{gen_grad_map} lattice
format. It is retained here for reference and to document the field convention used by lattices that
predate the curved-coordinate formalism (\sref{s:gen.grad.phys}). Old \vn{gen_grad_map} lattice
files can be converted to the current \vn{gen_gradients} format with the
\vn{gen_grad_map_to_gen_gradients} utility.

The electric and magnetic fields are both described by a scalar potential\footnote
  {
Notice the negative sign here and in \Eq{ppmpp} compared to Venturini et al.\cite{b:gen.grad}. This
is to keep the definition of the electric scalar potential $\psi_E$ consistent with the normal
definition.
  }
\begin{equation}
  \bfB = -\nabla \, \psi_B, \qquad \bfE = -\nabla \, \psi_E
  \label{bpep2}
\end{equation}
The scalar potential is then decomposed into azimuthal components
\begin{equation}
  \psi = \sum_{m = 1}^\infty \psi_{m,s} \, \sin(m \theta) + \sum_{m = 0}^\infty \psi_{m,c} \, \cos(m \theta)
\end{equation}
where the $\psi_{m,\alpha}$ ($\alpha = c,s$) are characterized by a using functions
$C_{m,\alpha}(z)$ which are functions along the longitudinal $z$-axis.
\begin{equation}
  \psi_{m,\alpha} = \sum_{n = 0}^\infty \frac{(-1)^{n+1} m!}{4^n \, n! \, (n+m)!} 
  \, \rho^{2n+m} \, C_{m,\alpha}^{[2n]}(z) 
  \label{ppmpp}
\end{equation}
The notation $[2n]$ indicates the $2n$\Th derivative of $C_{m,\alpha}(z)$.

From \Eq{ppmpp} the field is given by
\begin{align}
  B_\rho   &= \sum_{m = 1}^{\infty} \sum_{n = 0}^\infty \frac{(-1)^n \, m! \, (2n+m)}{4^n \, n! \, (n+m)!}
    \rho^{2n+m-1} \left[ C^{[2n]}_{m,s}(z) \, \sin m\theta + C^{[2n]}_{m,c}(z) \, \cos m\theta \right] + \CRNO
    & \hspace{25 em} \sum_{n = 1}^\infty \frac{(-1)^n \, 2n}{4^n n! \, n!} \rho^{2n-1} \, C^{[2n]}_{0,c}(z) \CRNO
  B_\theta &= \sum_{m = 1}^{\infty} \sum_{n = 0}^\infty \frac{(-1)^n \, m! \, m}{4^n \, n! \, (n+m)!}
    \rho^{2n+m-1} \left[ C^{[2n]}_{m,s}(z) \, \cos m\theta - C^{[2n]}_{m,c}(z) \, \sin m\theta \right] \\
  B_z      &= \sum_{m = 0}^{\infty} \sum_{n = 0}^\infty \frac{(-1)^n \, m!}{4^n \, n! \, (n+m)!}
    \rho^{2n+m} \left[ C^{[2n+1]}_{m,s}(z) \, \sin m\theta + C^{[2n+1]}_{m,c}(z) \, \cos m\theta \right] \nonumber
\end{align}
Even though the scalar potential only involves even derivatives of $C_{m,\alpha}$, the field is
dependent upon the odd derivatives as well. The multipole index $m$ is such that $m = 0$ is for
solenoidal fields, $m = 1$ is for dipole fields, $m = 2$ is for quadrupolar fields, etc. The
\vn{sin}--like generalized gradients represent normal (non-skew) fields and the \vn{cos}--like one
represent skew fields. The on-axis fields at $x=y=0$ are given by:
\begin{equation}
  (B_x, B_y, B_z) = (C_{1,c}, C_{1,s}, -C^{[1]}_{0,c})
\end{equation}

The magnetic vector potential for $m = 0$ is constructed such that only $A_\theta$ is non-zero
\begin{align}
  A_\rho   &= 0 \CRNO
  A_\theta &= \sum_{n=1}^\infty 
    \frac{(-1)^{n+1} \, 2n}{4^n \, n! \, n!} \rho^{2n-1} \, C_{0,c}^{[2n-1]} \\
  A_z      &= 0 \nonumber
\end{align}
For $m \ne 0$, the vector potential is chosen so that $A_\theta$ is zero.
\begin{align}
  A_\rho   &= \sum_{m = 1}^{\infty} \sum_{n=0}^\infty 
    \frac{(-1)^{n} \, (m-1)!}{4^n \, n! \, (n+m)!} \rho^{2n+m+1} \, 
    \left[ C_{m,s}^{[2n+1]} \cos(m \theta) - C_{m,c}^{[2n+1]} \, \sin(m \theta) \right] \CRNO 
  A_\theta &= 0 \\
  A_z      &= \sum_{m = 1}^{\infty} \sum_{n=0}^\infty 
    \frac{(-1)^n \, (m-1)! \, (2n+m)}{4^n \, n! \, (n+m)!} \rho^{2n+m} \, 
    \left[ -C_{m,s}^{[2n]} \cos(m \theta) + C_{m,c}^{[2n]} \, \sin(m \theta) \right]
  \nonumber
\end{align}

The functions $C_{m,\alpha}(z)$ are characterized by specifying $C_{m,\alpha}(z_i)$ and derivatives
at equally spaced points $z_i$, up to some maximum derivative order $N_{m,\alpha}$ chosen by the
user. Interpolation is done by constructing an interpolating polynomial (``non-smoothing spline'')
for each GG of order $2N_{m,\alpha}+1$ for each
interval $[z_i, z_{i+1}]$ which has the correct derivatives from $0$ to $N_{m,\alpha}$ at points $z_i$ and
$z_{i+1}$. The coefficients of the interpolating polynomial are easily calculated by inverting the
appropriate matrix equation.

The advantages of a generalized gradient map over a cylindrical or Cartesian map decomposition come
from the fact that with generalized gradients the field at some point $(x,y,z)$ is only dependent
upon the value of $C_{m,\alpha}(z)$ and derivatives at points $z_i$ and $z_{i+1}$ where $z$ is in
the interval $[z_i, z_{i+1}]$. This is in contrast to the cylindrical or Cartesian map decomposition
where the field at any point is dependent upon {\em all} of the terms that characterize the
field. This ``\vn{locality}'' property of generalized gradients means that calculating coefficients
is easier (the calculation of $C_{m,\alpha}(z)$ at $z_i$ can be done using only the field near $z_i$
independent of other regions) and it is easier to ensure that the field goes to zero at the
longitudinal ends of the element. Additionally, the evaluation is faster since only coefficients to
either side of the evaluation point contribute. The disadvantage of generalized gradients is that
since the derivatives are truncated at some order $N_{m,\alpha}$, the resulting field does not satisfy
Maxwell's equations with the error as a function of radius scaling with the power $\rho^{m+N_{m,\alpha}}$.

It is sometimes convenient to express the fields in terms of Cartesian coordinates. For sine like
even derivatives $C_{m,s}^{[2n]}$ the conversion is
\begin{align}
  \left( B_x, B_y \right) &= \left( \cos\theta \, B_\rho - \sin\theta \, B_\theta, \,
    \sin\theta \, B_\rho + \cos\theta \, B_\theta \right) \CRNO
  &= \frac{(-1)^n \, m!}{4^n n! \, (n+m)!} \, C_{m,s}^{[2n]} \, \Big[ (n+m) \, (x^2 + y^2)^n \,
    \left( S_{xy}(m-1), \, C_{xy}(m-1) \right) \, + 
    \label{bbtb} \\
  &\hspace{13em}  n \, (x^2 + y^2)^{n-1} \, 
    \left( S_{xy}(m+1), \, -C_{xy}(m+1) \right) \Big] \nonumber
\end{align}
and for the sine like odd derivatives $C_{m,s}^{[2n+1]}$
\begin{equation}
  B_z = \frac{(-1)^n \, m!}{4^n n! \, (n+m)!}
    (x^2 + y^2)^n \, C^{[2n+1]}_{m,s}(z) \, S_{xy}(m)
\end{equation}
where the last term in \Eq{bbtb} is only present for $n > 0$.
\begin{align}
  S_{xy}(m) &\equiv \rho^m \, \sin m\theta = 
    \sum_{r=0}^{2r \le m-1} (-1)^r \begin{pmatrix} m \cr 2r+1 \end{pmatrix} \,
    x^{m-2r-1} \, y^{2r+1} \CRNO
  C_{xy}(m) &\equiv \rho^m \, \cos m\theta = 
    \sum_{r=0}^{2r \le m} (-1)^r \begin{pmatrix} m \cr 2r \end{pmatrix} \,
    x^{m-2r} \, y^{2r} 
\end{align}
The conversion for the cosine like derivatives is:
\begin{align}
  \left( B_x, B_y \right) &= 
    \frac{(-1)^n \, m!}{4^n n! \, (n+m)!} \, C_{m,c}^{[2n]} \, \Big[ (n+m) \, (x^2 + y^2)^n \,
    \left( C_{xy}(m-1), \, -S_{xy}(m-1) \right) \, + \CRNO
  &\hspace{13em}  n \, (x^2 + y^2)^{n-1} \, 
    \left( C_{xy}(m+1), \, S_{xy}(m+1) \right) \Big] \\
  B_z &= \frac{(-1)^n \, m!}{4^n n! \, (n+m)!}
    (x^2 + y^2)^n \, C^{[2n+1]}_{m,c}(z) \, C_{xy}(m) \nonumber
\end{align}

%-----------------------------------------------------------------
